% ================================================================
%  ENCUENTRO 2 — Miércoles 29 de abril de 2026
% ================================================================
\chapter{Segundo Encuentro}

\begin{destacado}
    \textbf{Fecha:} Miércoles, 29 de abril de 2026 \quad
    \textbf{Horario:} 6:20 a.m.\ -- 8:00 a.m. \quad
    \textbf{Asistentes:} Thomas, Miguel, Kevin, Isabella, Esteban, Samuel Angulo
\end{destacado}

\medskip

Este segundo encuentro estuvo centrado en la apropiación inicial del telescopio como
instrumento de trabajo del semillero. La sesión combinó revisión de normas básicas de
uso, organización de consultas individuales, práctica de montaje y una primera salida
de observación terrestre. También permitió recoger intereses temáticos del grupo para
orientar las clases teóricas posteriores.

% ---------------------------------------------------------------
\section{Apertura de la sesión}
% ---------------------------------------------------------------

El encuentro inició en el salón de clases con una socialización de pautas básicas para
el manejo seguro y responsable del telescopio. Más que memorizar instrucciones, se
buscó que cada integrante comprendiera la razón técnica y preventiva detrás de cada
norma, de modo que el cuidado del instrumento se convierta en una práctica consciente.

\begin{concepto}{Aprender a usar un telescopio también es aprender a cuidarlo}{con:cuidado}
    En astronomía escolar, el dominio del instrumento no se limita a saber apuntar o
    enfocar. También implica conocer riesgos, prevenir daños y desarrollar hábitos de
    trabajo rigurosos. Un telescopio bien cuidado amplía las posibilidades de
    observación del semillero y garantiza continuidad en los procesos formativos.
\end{concepto}

% ---------------------------------------------------------------
\section{Consultas asignadas sobre normas de uso}
% ---------------------------------------------------------------

Con el fin de distribuir responsabilidades y promover la participación activa, se
asignó a cada estudiante una consulta breve relacionada con una norma específica de
uso del telescopio. Para el próximo encuentro, cada integrante deberá explicar:

\begin{itemize}[leftmargin=1.5em]
    \item por qué la instrucción asignada es importante;
    \item qué riesgo técnico o de seguridad se evita al cumplirla;
    \item qué podría ocurrir si se ignora.
\end{itemize}

\begin{actividad}{Consultas individuales para el próximo encuentro}{act:consultas}
    \begin{itemize}[leftmargin=1.5em]
        \item \textbf{Miguel:} no apuntar al Sol sin un filtro solar certificado.
        \item \textbf{Thomas:} aclimatación del telescopio antes de su uso.
        \item \textbf{Kevin:} evitar el contacto directo con los elementos ópticos.
        \item \textbf{Isabella:} cubrir el objetivo y el ocular cuando el equipo no
              esté en uso.
        \item \textbf{Esteban:} no forzar los mecanismos de ajuste del telescopio.
        \item \textbf{Samuel Angulo:} funcionamiento del sistema de aumento del
              telescopio \textit{Decouvir Maksutov-Cassegrain MC~60/70}.
    \end{itemize}
\end{actividad}

\begin{nota}{Criterio pedagógico de la actividad}{nota:criterio-consultas}
    Esta estrategia permite que las normas de uso no se presenten solo como una lista
    de prohibiciones, sino como conocimientos comprensibles y argumentados por los
    propios estudiantes. Así, la seguridad del instrumento se convierte en una tarea
    compartida por todo el grupo.
\end{nota}

% ---------------------------------------------------------------
\section{Montaje inicial del telescopio}
% ---------------------------------------------------------------

Luego de la socialización, se realizó el ensamblaje del telescopio dentro del salón.
Durante esta sesión, \textbf{Thomas} asumió la responsabilidad principal del montaje,
con el propósito de minimizar errores mientras el grupo se familiarizaba con el
procedimiento.

Aunque esta decisión permitió trabajar con mayor seguridad, también dejó clara la
necesidad de que todos los integrantes aprendan a montar y revisar el equipo de manera
autónoma.

\begin{idea}{Capacitación específica de ensamblaje y operación}{idea:capacitacion-montaje}
    Se propone programar una sesión práctica dedicada exclusivamente al montaje,
    desmontaje y operación básica del telescopio. Esta capacitación permitiría:
    \begin{itemize}[leftmargin=1.5em]
        \item garantizar que todos los integrantes dominen el procedimiento;
        \item establecer una verificación cruzada para reducir errores;
        \item fortalecer la autonomía del semillero ante futuras observaciones o ante
              la ausencia de algún integrante clave.
    \end{itemize}
    La fecha y la logística de esta capacitación quedan pendientes de definición.
\end{idea}

% ---------------------------------------------------------------
\section{Práctica de observación terrestre}
% ---------------------------------------------------------------

Una vez montado el equipo, el grupo se trasladó a la rampa del colegio por ofrecer un
campo visual amplio sobre la ciudad. Allí se realizaron pruebas de enfoque y se tomó
registro fotográfico de edificaciones y puntos de referencia lejanos.

Las condiciones atmosféricas no fueron favorables para una observación detallada:
el cielo permaneció muy nublado y la visibilidad fue limitada. Por ello, la práctica
se orientó principalmente a reconocer el comportamiento mecánico del telescopio,
ajustar el enfoque y explorar su uso en observaciones terrestres de corta complejidad.

\begin{concepto}{La observación terrestre como entrenamiento inicial}{con:observacion-terrestre}
    Antes de realizar observaciones astronómicas, es útil practicar con objetos
    terrestres lejanos. Esto ayuda a comprender el enfoque, la estabilidad del trípode,
    el campo visual y la orientación de la imagen, habilidades necesarias para luego
    localizar cuerpos celestes con mayor precisión.
\end{concepto}

% ---------------------------------------------------------------
\section{Cierre y exploración de intereses}
% ---------------------------------------------------------------

Al regresar al salón, se guardó el telescopio siguiendo los protocolos establecidos.
Como actividad de cierre, se propuso una pregunta abierta para identificar temas que
despiertan curiosidad en el grupo y que podrían convertirse en futuras clases,
consultas o proyectos de investigación:

\begin{destacado}
    \textbf{Pregunta orientadora:} \textit{¿Qué tema de astronomía han escuchado,
    pero aún no comprenden bien, y les gustaría aprender en el semillero?}
\end{destacado}

Las respuestas registradas muestran intereses diversos que abarcan tanto astronomía
como exploración espacial y cosmología:

\begin{itemize}[leftmargin=1.5em]
    \item agujeros de gusano;
    \item naves espaciales;
    \item misión Artemis II y otras misiones cercanas;
    \item astrobiología y posibilidad de vida extraterrestre;
    \item formación de galaxias;
    \item Gran Atractor.
\end{itemize}

\begin{nota}{Lectura pedagógica de las respuestas}{nota:intereses}
    Los temas elegidos revelan que el grupo siente atracción tanto por preguntas
    científicas profundas como por asuntos de alta presencia mediática. Esto ofrece
    una oportunidad valiosa: partir de la curiosidad espontánea de los estudiantes
    para introducir conceptos rigurosos de física, astronomía y método científico.
\end{nota}

% ---------------------------------------------------------------
\section{Proyección del próximo encuentro}
% ---------------------------------------------------------------

Como resultado de lo trabajado en esta sesión, se acordó orientar el siguiente
encuentro hacia una clase teórica introductoria a cargo del profesor Daniel, con énfasis en:

\begin{itemize}[leftmargin=1.5em]
    \item nuestra posición en el universo;
    \item escalas de tamaño y distancias cósmicas;
    \item bases generales para comprender el lugar de la Tierra, el sistema solar,
          la galaxia y el universo observable.
\end{itemize}

\begin{compromiso}{Compromisos para la siguiente sesión}{comp:clase2}
    \begin{itemize}[leftmargin=1.5em]
        \item Cada estudiante preparará la consulta asignada sobre normas de uso del
              telescopio.
        \item Se continuará fortaleciendo el manejo técnico del instrumento desde una
              perspectiva colectiva y responsable.
        \item Se empezará a organizar una ruta temática que responda a los intereses
              expresados por el grupo en esta sesión.
    \end{itemize}
\end{compromiso}